% !PRINTED 393
\weberpage{393}
Accordingly we have:

4. \gesperrt{If $\frac14 D \equiv 1\ (\mathrm{mod}\ 8)$, then every odd number is a norm residue of 4 and of 8.}

\gesperrt{If $\frac14 D \equiv 0\ (\mathrm{mod}\ 8)$, then only the numbers of the form $4n+1$ are norm residues of 4, and the numbers of the form $8n+1$ norm residues of 8.}

\gesperrt{If $\frac14 D \equiv 2\ (\mathrm{mod}\ 4)$, then all odd numbers are norm residues of 4.}

\gesperrt{If $\frac14 D \equiv 2$, or $\equiv 6\ (\mathrm{mod}\ 8)$, then in the first case the numbers $8n+1$, $8n-1$, in the second case the numbers $8n+1$, $8n+3$ are norm residues of 8.}

\gesperrt{If $\frac14 D \equiv 3, 4, 7\ (\mathrm{mod}\ 8)$, then all and only the numbers of the form $4n+1$ are norm residues of 4 and of 8.}

5. \gesperrt{If $p$ is an odd prime, and an $a$ not divisible by $p$ is a norm residue of $p^k$ (for any positive exponent $k$), then $a$ is also a norm residue of $p^{k+1}$.}

For the proof let

\[
N(\omega) = a + p^k b,
\]
\[
\omega_1 = \omega\,(1 + x\,p^k),
\]

where $x$ is rational. Consequently

\[
\begin{aligned}
N(\omega_1)&=(a+p^k b)(1+xp^k)^2,\\
&\equiv a+p^k(b+ax)\ (\mathrm{mod}\ p^{k+1}),
\end{aligned}
\]

and if $x$ is determined from the congruence $b + a x \equiv 0\ (\mathrm{mod}\ p)$, then it follows that

\[
N(\omega_1) \equiv a\ (\mathrm{mod}\ p^{k+1}).
\]

6. \gesperrt{If $D$ is odd, the same still holds for $p = 2$ as well.}

For in this case every odd $a$ is a norm residue of 8. We therefore put

\[
N(\omega) = a + 2^k b, \qquad k \geqq 3,
\]
\[
\omega_1 = \omega \left(1 + 2^k x\,\frac{1+\sqrt{D}}{2}\right),
\]
\[
N(\omega_1) \equiv (a + 2^k b)(1 + 2^k x)\ (\mathrm{mod}\ 2^{k+1}),
\]

and if one assumes $a x + b \equiv 0\ (\mathrm{mod}\ 2)$, then it follows that

\[
N(\omega_1) \equiv a\ (\mathrm{mod}\ 2^{k+1}).
\]

7. \gesperrt{If $D$ is even and $a$ a norm residue of $2^k \geqq 8$, then $a$ is also a norm residue of $2^{k+1}$.}
